\chapter{Model and Simulation Design}\label{ch:model}
%% FUNCTION: Explains how we measure how often round robin achieves EFX, in the order in which the
%% simulation itself proceeds: a scenario (agents and goods), values for every agent and good, the
%% round robin allocation and its evaluation, and finally the design grid with the reference cell.
%% Random allocation and other picking sequences are introduced in Chapter 5, not here.

Our simulation proceeds in four steps, and this chapter follows them. A scenario fixes how many
agents share how many goods (Section~\ref{sec:scenario}). For every agent and good, we draw a
value (Section~\ref{sec:instances}). Round Robin allocates the goods, and we measure how close the
allocation comes to EFX (Section~\ref{sec:evaluation}). Repeating these steps many times for every
combination of parameters gives the design (Section~\ref{sec:design}).

\section{Scenario}\label{sec:scenario}

A scenario consists of $n$ agents and $m$ goods with $m>n$; for $m\leq n$, every allocation by
Round Robin is trivially EFX (Section~\ref{sec:notions}). Following
\textcite[p.~16]{manurangsi2021}, we write $m=rn+q$ with $r=\lfloor m/n\rfloor$ full rounds and a
remainder $0\leq q<n$. Under Round Robin, agents $1,\dots,q$ receive $r+1$ goods and the others
$r$ goods. For $q=0$, the number of goods is divisible by the number of agents, and every agent
receives the same number of goods.

The four quantities are linked by $m=rn+q$, so only some comparisons separate one effect from
another. Comparing different numbers of goods at a fixed number of agents changes $r$ or $q$ or
both, and comparing different numbers of agents at a fixed number of goods changes all of them.
We therefore never speak of an effect of $m$ alone. Instead, we compare
\begin{itemize}
  \item different $r$ at fixed $q$ (Section~\ref{sec:goods}),
  \item different $q$ at fixed $r$ (Section~\ref{sec:divisibility}), and
  \item different $n$ at fixed $r$ and $q$, that is, more agents with the same number of goods
  per agent (Section~\ref{sec:agents}).
\end{itemize}

\section{Values}\label{sec:instances}

We generate the value that an agent attributes to a good in two steps. In the first step, we
determine how the agents rank the goods. For every agent $i$ and good $g$, we draw a score $\Zig$
that consists of a common component $\Wg$, shared by all agents, and a private component $\etaig$,
which reflects the individual taste of agent $i$:
\[
  \Zig=\sqrt{\rho}\,\Wg+\sqrt{1-\rho}\,\etaig,\qquad \Wg,\etaig\sim\mathcal{N}(0,1).
\]
All $\Wg$ and $\etaig$ are mutually independent, and the parameter $\rho\in[0,1]$ sets the weight
of the common component. Two properties motivate this construction: $\Var(\Zig)=1$ for every
$\rho$, and $\Corr(\Zig,Z_{jg})=\rho$ for $i\neq j$. Thus, $\rho$ only influences how similarly the
agents rank the goods, not the distribution of a single score. For $\rho=0$, the rankings of
different agents are independent and uniformly distributed over all orderings of the goods,
whereas $\rho=1$ makes all rankings identical. These two endpoints coincide with the models
\textit{full independence} and \textit{full correlation} of \textcite[p.~74]{bouveret2011}. We call
$\rho$ the \textit{agreement}.

In the second step, we stretch $\Zig$ to obtain the value
\[
  v_i(g)=\exp(\theta\,\Zig),
\]
where $\theta>0$ controls how strongly the values are spread; we call it the
\textit{dispersion}. Each $v_i(g)$ is lognormally distributed, and this distribution does not
depend on $\rho$. The values are strictly positive and continuously distributed, so ties occur
with probability zero. As $\theta\to0$, all values tend to $1$, and every good is worth almost the
same. We chose the lognormal distribution because its values are positive and a single parameter
controls their spread. This is our own modelling choice. Similar models that combine a common and
an individual component appear in \textcite[p.~1406]{dickerson2014} and
\textcite[p.~8]{branzei2025}; in the latter, the parameter is a mixing weight rather than a
correlation.

We report both parameters on interpretable scales. Agreement is reported as the Spearman rank
correlation between the values of two agents over the goods,
\[
  \rho_S=\frac{6}{\pi}\arcsin\Big(\frac{\rho}{2}\Big),
\]
and dispersion as the Gini coefficient of the values of a single agent over the goods,
\[
  G=2\,\Phi\big(\theta/\sqrt{2}\big)-1,
\]
where $\Phi$ denotes the standard normal distribution function. Each scale reflects exactly one
parameter. Table~\ref{tab:axes} lists the levels used in this thesis. We report every Gini value
together with $\theta$ in parentheses.

\begin{table}[ht]
  \centering
  \caption{Parameter levels and their interpretable counterparts.}\label{tab:axes}
  \begin{tabular}{lcccccc}
    \toprule
    Agreement $\rho$ & 0 & 0.25 & 0.5 & 0.75 & 1 & \\
    Spearman $\rho_S$ & 0 & 0.24 & 0.48 & 0.73 & 1 & \\
    \midrule
    Dispersion $\theta$ & 0.1 & 0.25 & 0.5 & 0.75 & 1.0 & 1.5 \\
    Gini $G$ & 0.06 & 0.14 & 0.28 & 0.40 & 0.52 & 0.71 \\
    \bottomrule
  \end{tabular}
\end{table}

The Spearman scale describes the distribution, not a single instance. With a finite number of
goods, the rank correlation between two agents in a single instance can differ considerably from
$\rho_S$. We therefore read a Spearman value as \textit{the values are drawn from a world with
rank correlation $\rho_S$}, and not as \textit{the agents agree to $\rho_S$}. All measured results
rest on this one family of distributions; Section~\ref{sec:limitations} discusses this limitation.

\section{Allocation and Evaluation}\label{sec:evaluation}

Round Robin receives only the values and returns an allocation (Section~\ref{sec:setting}). It
does not know $\rho$, $\theta$ or the cell of the design from which an instance comes. Since it
only compares values, it depends on the instance only through the rankings. Because
$x\mapsto\exp(\theta x)$ is strictly increasing, $\theta$ does not change any ranking. Hence the
Round Robin allocation of an instance is the same for every level of dispersion (proven):
dispersion changes how a fixed allocation is evaluated, never the allocation itself. Ties occur
with probability zero, so the tie-breaking rule plays no role, and since the agents are
exchangeable in our model, letting agent $1$ pick first is no restriction.

For every allocation, we compute the ratio $\aij$ for every pair of agents and the target
quantity $\astar$ from Section~\ref{sec:notions}. By construction, $\astar=1$ exactly when the
allocation is EFX, and $\astar<1$ measures how far it misses EFX. A cell
$X=(n,m,\rho,\theta)$ is the unit of our simulation. For each cell, we generate $R$ instances
$k=1,\dots,R$ with values $\astar_k$ and describe them by the survival function
\[
  \Surv(\alpha)=\Pr\big[\astar\geq\alpha \mid X\big],\qquad \alpha\in(0,1],
\]
and its empirical estimator
\[
  \Survhat(\alpha)=\frac{1}{R}\sum_{k=1}^{R}\mathbf{1}\{\astar_k\geq\alpha\}.
\]
$\Survhat(\alpha)$ is the share of instances that come at least as close to EFX as $\alpha$. It is
non-increasing in $\alpha$, and its value at $\alpha=1$ is the share of instances that are EFX.
The curve below $\alpha=1$ shows how far the remaining instances are from EFX.

We summarise each cell by two key figures. The first is the \textit{EFX rate} $\Survhat(1)$. The
second is the $5\%$ quantile of $\astar$ among the instances with $\astar<1$, which we call the
\textit{conditional $5\%$ quantile}. It answers the question: if Round Robin fails, how far from
EFX are the worst five percent of the failures? We report it only if a cell contains at least
$200$ failures and show cells with fewer failures in grey. We do not use the mean of $\astar$ as a
key figure, because it mixes how often an allocation is EFX with how far it misses EFX. Since a
single quantile shows only one point of the failures, our results always read the whole survival
curve as well.

\section{Design and Reference Cell}\label{sec:design}

The design combines $n\in\{2,3,4,5\}$ agents with $m=n+1,\dots,4n$ goods, the five levels of
agreement and the six levels of dispersion from Table~\ref{tab:axes}. This gives $42$ pairs
$(n,m)$ and $1{,}260$ cells, each with $R=20{,}000$ instances. With this number of instances, the
standard error $\sqrt{\Survhat(1)(1-\Survhat(1))/R}$ of every EFX rate is at most $0.0035$,
attained at an EFX rate of $0.5$. The same instances are also allocated by the mechanisms of
Chapter~\ref{ch:random}.

All results refer to a common \textit{reference cell}: $n=4$ agents and $m=12$ goods, so $r=3$
and $q=0$, with Spearman $0.48$ ($\rho=0.5$) and Gini $0.52$ ($\theta=1$). Each section of
Chapter~\ref{ch:results} varies one of these quantities and keeps the others fixed. The reference
cell lies in the region where the existence of an EFX allocation is open: four agents, at least
four goods more than agents, and less than full agreement (Section~\ref{sec:existence}). There,
and only there, our interpretation rests on the assumption that an EFX allocation exists.

The seed of an instance depends only on a master seed and on $k$. Instance $k$ therefore uses the
same random numbers $\Wg$ and $\etaig$ for every level of agreement and dispersion and for every
mechanism (common random numbers). Curves over agreement, dispersion and mechanisms are thus
paired, not independent. Differences between them are estimated more precisely, but counts over
many cells are not independent confirmations. Over dispersion, the pairing is perfect: the
allocation stays the same, and only its evaluation changes.

We ran the simulation twice. A first, exploratory run with $R=2{,}000$ and seed $12345$ served
only to develop hypotheses. All numbers reported in this thesis come from a separate confirmatory
run with $R=20{,}000$ and seed $20260927$, which we fixed before the run. Where we fixed a counting
rule only after having seen individual cells, we mark the result as exploratory. The simulation
stores one file per cell with the raw matrix of $\aij$ values, so every number in this thesis can
be recomputed from the stored cells.
% TODO: public repository, Python and NumPy versions (NumPy does not guarantee identical random numbers across versions)
